% !TEX root =  ../Dissertation.tex

\chapter{Discussion}

\section{Limitations and Lessons}
There are an endless number of editions and changes that the web application would have received if there had been no time constraints. That being said, a few in particular should be mentioned that would have had a significant impact.
Firstly, if development had been finished sooner, users could have used the application longer, giving more reliable data. Not only would the data have given readers a better idea about the long-term effect of these features, but users would have had more time to give feedback that could have been considered when making another round of changes. 
One of the changes immediately requested that should have been taken into account was to ensure the application could be used on the phone. When the web application was first sent out, if it was loaded on a smaller screen, it would only display the keyboard, and the entire interface would be unusable. This massive issue was fixed on the second day by adjusting the CSS for smaller screens. 
There were also accessibility issues when users used the web application on the phone. One potential participant had to drop out of the experiment as they could only use their phone during the week, and their visual impairment meant that they used large text when using their phone, which is a setting within many smartphones. This caused the text to appear fragmented on the page, and unfortunately, it couldn’t be fixed before the experiment ended. If development were to start again from scratch, it would have been approached with more mindfulness to accessibility for all users.
 Perhaps naively, the web application was only made with a desktop in mind, and as development had taken place on a 1440p screen, there was an assumption that it would still work on smaller devices. The issue was so prevalent because of all the pageviews on the site, 42 percent of which were visited on the phone. The fix that made the application usable on the phone also did not have the neatest layout and required scrolling to get to the submit word button, particularly on eight-letter wordles, where a user receives nine guesses. In the future, when making a game like this again, it would be more important that building the web application from the start with smaller screens in mind would lead to higher satisfaction with users.
Leaderboards also underwent an iteration when users started playing. Initially, only one leaderboard contained the data for all word lengths. It was quickly realised that this discouraged users from playing some of the longer word lengths, as these are harder to get in a small number of guesses. This was fixed quickly, but this may have discouraged people if they were not often checking the leaderboards. 
Another issue with the Leaderboards is that the tiebreaker for people with the same average score is their average time to complete a game. The average time is calculated by checking the end time of each game minus the start time of each game in seconds. Then, it was used as their average time. While this seemed logical at the time, something that hadn’t been considered was when a user leaves a game unfinished overnight. This led to very long average times that did not indicate how long a user takes to complete a game. 
Another issue that needed to be fixed by the end of development was invalidating guesses. A user can type any assortment of letters as a guess, and it will accept it as a word. This is because many word lists were tried during development but would not take some words as valid during testing. Some libraries needed the full range of answers you can get in the game. This led to issues where people had to ask for their game to be deleted from the database as they knew the correct word but could not use it as an answer. After many different libraries were tested, it was decided that it would be best to avoid having this guess checker to stop people from getting frustrated with actual words that are not considered valid.
 While this is a limitation to achieving parity with the original Wordle game, it led to users' creativity when trying to complete the more complex challenges that higher-length words provided. Users would throw away a guess to try and find where the vowels were in the eight-letter words, making it easier for them. This development of a ‘meta’ of how to complete higher-length wordles was an exciting example of the innovation users can find when the developer allows it. 
 
\section{Future Works}
Moving forward with the application, many additions could be made to improve it, making the experience more enjoyable. It would be interesting to develop new games from the application and make it less inspired by Wordle and instead NYTGames. NYTGames acquired Wordle a few months after its creation and added it to a similar set of verbal and non-verbal reasoning games. This includes Crosswords, Mini Crosswords, Wordle, a game called connections where you have to find the links between 16 words, a more complex word search and many others. Adding these many features would likely increase retention and allow users to pick which games they enjoy the most and only play those. In the original scope of this project, it was intended to have two more games, and it would be ideal to develop a few more so users can play Wordle. This would also allow us to know how much of a factor it was of people needing to be retained because they were bored of Wordle. 

To make the experiment have more compelling data with more usability, more retention techniques could be developed, and when an account is created, the user is set 3 retention techniques from a list of 10. With this, you could make more precise decisions about which strategies were most effective, and creators of web applications in the daily word game genre could make more accurate decisions on how to spend development time creating retention techniques that will retain users. 

\section{Threats to Validity}
While the data tracked on both websites indicate that the study has shown the hypothesis to be true based on user satisfaction and data metrics. It should be apparent that only tracking user retention for seven days could lead to misunderstandings about how well the web application would perform in the long term. From the first seven days, we can be fairly certain that The Gauntlet would outperform Wordle with no login, but how much and whether the Gauntlet would also lose its player base remains unclear. 
The small sample size of users may also raise questions about how well the website would have retained players if the web application had started with 1000 or more users. From the data of other games, it is clear that user retention would not remain steady, as it is expected that after thirty days, user retention will usually drop another 3-4 percent from the initial user base.\cite{mistplay2024}

\section{Conclusion}
The goal of developing both web applications was to explore how the implementation of retention strategies impacts user satisfaction, engagement, and retention. The data metrics and user surveys show that incorporating retention techniques into The Gauntlet led to significantly better performance compared to Wordle without a login system, particularly in terms of user satisfaction, engagement time, and user retention. Features such as leaderboards, achievements, and enhanced functionality played a key role in driving these improvements. However, there is still room for further exploration. Future work could focus on developing adaptive retention strategies tailored to individual users. This approach would allow developers to identify the most effective techniques for maximizing user satisfaction and retention on a more personalized level.